\begin{python0}
from solutions import *; clear()
\end{python0}
\sectionthree{The File Cursor/Pointer}

Recall that when you were reading the file

\begin{console}
hello world
0 3.14
42
\end{console}

with this program

\begin{console}
#include <iostream>
#include <fstream>

int main()
{   
    std::ofstream f("hw.txt", std::ios::app);

    f << "hello world" << std::endl;
    int x = 0;
    double y = 3.14;
    char z[] = "c++ rules";
    f << x << ' ' << y << std::endl;
    f << z << std::endl;

    f.close();

    return 0;
}
\end{console}

The reading of data for each variable stops at a whitespace (space, tab, newline). The next read will continue where it previous stopped.

This tells you that each file remembers where the previous read stopped. This is called a file cursor.

When you open a file, the cursor points to the first character in the file. As you read, the cursor moves forward.

Let's look at the above file:

\begin{console}
hello world
0 3.14
42 
\end{console}

It looks like there are three "lines" in this file. Actually a file is made up of a bunch of bytes (or characters if you like). On the harddrive (or some other storage device), the above data might look like this on the device:
\begin{python}
from latextool_basic import *
p = Plot()
s = array(x0=0.5, y0=0.5, width=0.6, height=1.2, xs=['h', 'e', 'l', 'l', 'o', ' ', 'w', 'o', 'r', 'l', 'd', r'\textbackslash n', 0, ' ', 3, '.', 1, 4, r'\textbackslash  n', 4, 2, ' ',' ',' ',' ',' ',' ',' ',' '], celllinewidth=0.05)
p.add(s)
p += Rect(13.1, 0.51, 13.7, 1.69, background='black')
#p+= Grid(x0=-1, y0=-2, x1=15, y1=5)
#p += Line(points=[(0.85,-1),(0.85,0.5)], linewidth=0.3, endstyle='>')
print(p)
\end{python}

(It's actually a lot more complex than this but this is enough.)

Each square in the above is a "slot" in your storage device to store a byte (or character) of data. You can think of your storage device as a pretty huge array if you like. The black square denotes the "end of the file" -- EOF. Initially, when the file is opened, its file cursor points to the first byte, i.e.the character h:

\begin{python}
from latextool_basic import *
p = Plot()
s = array(x0=0.5, y0=0.5, width=0.6, height=1.2, xs=['h', 'e', 'l', 'l', 'o', ' ', 'w', 'o', 'r', 'l', 'd', r'\textbackslash n', 0, ' ', 3, '.', 1, 4, r'\textbackslash  n', 4, 2, ' ',' ',' ',' ',' ',' ',' ',' '], celllinewidth=0.05)
p.add(s)
p += Rect(13.1, 0.51, 13.7, 1.69, background='black')
#p+= Grid(x0=-1, y0=-2, x1=15, y1=5)
p += Line(points=[(0.85,-1),(0.85,0.5)], linewidth=0.3, endstyle='>')
print(p)
\end{python}
If you execute this:

\tab[3em]{\verb!char x[100];!}\\
\tab[3em]{\verb!f >> x;!}\\

the characters starting from the character where the file cursor points to is read and put into array \texttt{x} until a whitespace is reached:

\begin{python}
from latextool_basic import *
p = Plot()
s = array(x0=0.5, y0=0.5, width=0.6, height=1.2, xs=['h', 'e', 'l', 'l', 'o', ' ', 'w', 'o', 'r', 'l', 'd', r'\textbackslash n', 0, ' ', 3, '.', 1, 4, r'\textbackslash  n', 4, 2, ' ',' ',' ',' ',' ',' ',' ',' '], celllinewidth=0.05)
p.add(s)
p += Rect(13.1, 0.51, 13.7, 1.69, background='black')
#p+= Grid(x0=-1, y0=-2, x1=15, y1=5)
p += Line(points=[(3.85,-1),(3.85,0.5)], linewidth=0.3, endstyle='>')
print(p)
\end{python}
If you do this again:

\tab[3em]{\texttt{f >> x;}}\\

the file cursor will read and ignore whitespaces (spaces, tabs,
newlines) until it sees a non-whitespace

\begin{python}
from latextool_basic import *
p = Plot()
s = array(x0=0.5, y0=0.5, width=0.6, height=1.2, xs=['h', 'e', 'l', 'l', 'o', ' ', 'w', 'o', 'r', 'l', 'd', r'\textbackslash n', 0, ' ', 3, '.', 1, 4, r'\textbackslash  n', 4, 2, ' ',' ',' ',' ',' ',' ',' ',' '], celllinewidth=0.05)
p.add(s)
p += Rect(13.1, 0.51, 13.7, 1.69, background='black')
#p+= Grid(x0=-1, y0=-2, x1=15, y1=5)
p += Line(points=[(4.45,-1),(4.45,0.5)], linewidth=0.3, endstyle='>')
print(p)
\end{python}
and then read the data, putting the characters into \texttt{x}, until it sees a whitespace:

\begin{python}
from latextool_basic import *
p = Plot()
s = array(x0=0.5, y0=0.5, width=0.6, height=1.2, xs=['h', 'e', 'l', 'l', 'o', ' ', 'w', 'o', 'r', 'l', 'd', r'\textbackslash n', 0, ' ', 3, '.', 1, 4, r'\textbackslash  n', 4, 2, ' ',' ',' ',' ',' ',' ',' ',' '], celllinewidth=0.05)
p.add(s)
p += Rect(13.1, 0.51, 13.7, 1.69, background='black')
#p+= Grid(x0=-1, y0=-2, x1=15, y1=5)
p += Line(points=[(7.45,-1),(7.45,0.5)], linewidth=0.3, endstyle='>')
print(p)
\end{python}
The next time you read the file and put the data into an integer variable, the integer 0 will be placed in the variable, and the pointer moves to this point:

\begin{python}
from latextool_basic import *
p = Plot()
s = array(x0=0.5, y0=0.5, width=0.6, height=1.2, xs=['h', 'e', 'l', 'l', 'o', ' ', 'w', 'o', 'r', 'l', 'd', r'\textbackslash n', 0, ' ', 3, '.', 1, 4, r'\textbackslash  n', 4, 2, ' ',' ',' ',' ',' ',' ',' ',' '], celllinewidth=0.05)
p.add(s)
p += Rect(13.1, 0.51, 13.7, 1.69, background='black')
#p+= Grid(x0=-1, y0=-2, x1=15, y1=5)
p += Line(points=[(8.65,-1),(8.65,0.5)], linewidth=0.3, endstyle='>')
print(p)
\end{python}

And after reading a \texttt{double} into a \texttt{double} variable
it's here:

\begin{python}
from latextool_basic import *
p = Plot()
s = array(x0=0.5, y0=0.5, width=0.6, height=1.2, xs=['h', 'e', 'l', 'l', 'o', ' ', 'w', 'o', 'r', 'l', 'd', r'\textbackslash n', 0, ' ', 3, '.', 1, 4, r'\textbackslash  n', 4, 2, ' ',' ',' ',' ',' ',' ',' ',' '], celllinewidth=0.05)
p.add(s)
p += Rect(13.1, 0.51, 13.7, 1.69, background='black')
#p+= Grid(x0=-1, y0=-2, x1=15, y1=5)
p += Line(points=[(12.25,-1),(12.25,0.5)], linewidth=0.3, endstyle='>')
print(p)
\end{python}
And finally after reading an \texttt{int} value it's here:

\begin{python}
from latextool_basic import *
p = Plot()
s = array(x0=0.5, y0=0.5, width=0.6, height=1.2, xs=['h', 'e', 'l', 'l', 'o', ' ', 'w', 'o', 'r', 'l', 'd', r'\textbackslash n', 0, ' ', 3, '.', 1, 4, r'\textbackslash  n', 4, 2, ' ',' ',' ',' ',' ',' ',' ',' '], celllinewidth=0.05)
p.add(s)
p += Rect(13.1, 0.51, 13.7, 1.69, background='black')
#p+= Grid(x0=-1, y0=-2, x1=15, y1=5)
p += Line(points=[(13.45,-1),(13.45,0.5)], linewidth=0.3, endstyle='>')
print(p)
\end{python}

The file cursor has reached the end-of-file.

